\section{Proof of Properties~\texorpdfstring{\ref{1}--\ref{4}}{1--4}}\label{sec:proof}
We remind the reader that $\Bbbk$ is an arbitrary commutative
ring. Now we are ready to start the inductive proof of the four
properties. For each property, we need to assume an earlier instance
of one or more properties. We begin with Property 1.


\begin{prop}\label{prop:extend}
Assume Property~\ref{1} for $(n-1,r)$. Then:
\begin{enumerate}[label=(\alph*)]
\item\label{prop6.1_a} For any $i$, $j$ with $1 \le i,j \le n$ and any given $\bB$ in
 $\E_\Bbbk(n,r-1)^i_j$, there exists some $\bA$ in
 $\E_\Bbbk(n,r)^i_j$ such that $\rho(\bA) = \bB$.

\item\label{prop6.1_b} 
 \begin{enumerate}[label=(\roman*)]
 \item\label{prop6.1_bi} Fix some $j$ in $\{1, \dots, n\}$. Suppose given
 the data
 \[\bB(1), \dots, \bB(n)\] with $\bB(i)$ in
 $\E_\Bbbk(n,r-1)^i_j$ for $i = 1, \dots, n$. Then there exist
 corresponding $\bA(i)$ in $\E_\Bbbk(n,r)^i_j$ for $i = 1, \dots, n$
 such that $\rho(\bA(i)) = \bB(i)$ for all $i$. 

 \item\label{prop6.1_bii} Similarly, fix some $i$ in $\{1, \dots, n\}$. Suppose
 given the data
 \[\bB(1), \dots, \bB(n)\] with 
 $\bB(j)$ in
 $\E_\Bbbk(n,r-1)^i_j$ for $j = 1, \dots, n$. Then there exist
 corresponding $\bA(j)$ in $\E_\Bbbk(n,r)^i_j$ for $j = 1, \dots, n$
 such that $\rho(\bA(j)) = \bB(j)$ for all $j$. 

 \item\label{prop6.1_biii} In either case~\ref{prop6.1_bi} or~\ref{prop6.1_bii}, the sum $\bA(1)
 + \cdots + \bA(n)$ extends $\bB(1) + \cdots + \bB(n)$.
 \end{enumerate}
 
\item\label{prop6.1_c} If in addition Property~\ref{2} holds for $(n,r-1)$ then
 Property~\ref{1} holds for $(n,r)$.
\end{enumerate}
\end{prop}

\begin{proof}
\ref{prop6.1_a}~This follows from Proposition~\ref{prop:n-reduce-iso}, which
 reduces the question to the problem of extending from
 $\E_\Bbbk(n-1,r-1)$ to $\E_\Bbbk(n-1,r)$, which is solved by the
 hypothesis.

\ref{prop6.1_b}~is immediate from part~\ref{prop6.1_a} and the linearity of $\rho$.

\ref{prop6.1_c}~Let $\bB$ in $\E_\Bbbk(n,r-1)$ be given. By the decomposition
 property for $(n,r-1)$, we can find $\bB(j)$ in
 $\E_\Bbbk(n,r-1)^n_j$ for $j = 1, \dots, n$ such that $\bB = \bB(1)
 + \cdots + \bB(n)$. By part~\ref{prop6.1_b}\ref{prop6.1_bii}, there exist corresponding
 $\bA(j)$ in $\E_\Bbbk(n,r)^n_j$ such that $\rho(\bA(j)) = \bB(j)$
 for all $j = 1, \dots, n$. Put $\bA = \bA(1) + \cdots +
 \bA(n)$. Then by linearity of $\rho$ it follows that $\rho(\bA) =
 \bB$, as required. This shows that $\bA$ is a last block row
 extension of $\bB$. The proof for any other block row or column is
 similar.
\end{proof}

Having dealt with the existence of extensions, we now consider the
question of their uniqueness.

\begin{lemm}\label{lem:rho-injects}
 Suppose that $n \le r$. Then
 \begin{enumerate}[label=(\alph*)]
 \item\label{lemma6.2_a} Extensions \textup{(}if any\textup{)} from $\E_\Bbbk(n,r-1)$
 to $\E_\Bbbk(n,r)$ are unique.
 \item\label{lemma6.2_b} Restriction $\rho: \E_\Bbbk(n,r) \to \E_\Bbbk(n,r-1)$ is
 injective.
 \item\label{lemma6.2_c} $\bzero$ is the only extension in $\E_\Bbbk(n,r)$ of $\bzero$.
 \end{enumerate}
\end{lemm}

\begin{proof}
\ref{lemma6.2_a}~Suppose that $\bA$ in $\E_\Bbbk(n,r)$ extends $\bB$ in
 $\E_\Bbbk(n,r-1)$. We use Proposition~\ref{prop:inv-properties}. Let
 $\bi = i_1 \cdots i_r$, $\bj = j_1 \cdots j_r$ be in $I(n,r)$. If
 $\vt(\bi) \ne \vt(\bj)$ then by Proposition~\ref{prop:inv-properties}\ref{prop3.4_c}, $a^\bi_\bj = 0$.

\goodbreak
 So assume for the rest of the proof that $\vt(\bi) = \vt(\bj)$. If
 $\bi$ has a repeated value, then by Proposition~\ref{prop:inv-properties}\ref{prop3.4_d}, $a^\bi_\bj = b^\bp_\bq$, where $\bp,
 \bq$ are obtained from $\bi, \bj$ by removing one of the duplicate
 values (from the same place). So such entries of $\bA$ are
 determined by $\bB$. If $r>n$, then all multi-indices in $I(n,r)$
 must have at least one duplicate value, so we are done in that case.

 We are left with the case $r=n$ and $\#(\bi) = n$ ($\bi$ has no
 repeated values). Then the same is true of $\bj$, and by Proposition~\ref{prop:inv-properties}\ref{prop3.4_b}, with $\bp = i_1 \cdots i_{n-1}$ and
 $\bq = j_1 \cdots j_{n-1}$ we have
 \[
 \textstyle \sum a^{\bp\,*}_{\bq\,j_r} = b^\bp_\bq .
 \]
 Exactly $n-1$ of the possible $n$ values from $\{1, \dots, n\}$
 appear in $\bi$; similarly for $\bj$. Hence at most one term in the
 above sum can be non-zero, because of value-type, so $a^\bi_\bj =
 b^\bp_\bq$. So in this case, $\bA$ is also determined by $\bB$. This
 proves part~\ref{lemma6.2_a}.

\ref{lemma6.2_b}~This follows from part~\ref{lemma6.2_a}. If $\bA \in \ker \rho$, then $\bA =
 \bzero$ by uniqueness.

\ref{lemma6.2_c}~This follows from part~\ref{lemma6.2_b}.
\end{proof}

\begin{rema}\label{rmk:sharper}
 Suppose that $n \le r+1$. If $\bC$ is any invariant in
 $\E_\Bbbk(n,r)$ with zero last block row, then $\bC = \bzero$. The
 same holds for any other block row or column. This follows from the
 proof of Lemma~\ref{lem:rho-injects}\ref{lemma6.2_a}; the assumption of a zero
 block row removes one degree of freedom from column slice sums, so
 the uniqueness conclusion is still valid, so $\bC=\bzero$.
\end{rema}

The next result gives conditions under which we can construct
extensions with prescribed partial information in a chosen block row
or column. This is a crucial technical result needed to prove
Property~\ref{2}. The chosen block row or column is controlled by its
free-pattern $F'(n,r) \cong D(n,r-1)$. We say that the prescribed
information is \emph{compatible} with the extension problem for a
given $\bB$ in $\E_\Bbbk(n,r-1)$ if there is some partial assignment
from a subset of some $F'(n,r)$ to $\Bbbk$ which gives the prescribed
information.



\begin{lemm}\label{lem:partial}
Assume Property~\ref{1} for $(n-1,r)$ and Property~\ref{4} for
$(n,r-1)$. Suppose that $\bB$ in $\E_\Bbbk(n,r-1)$ is given. Fix a
choice of block row \textup{(}respectively, block column\textup{)} and
a choice of any number of compatible prescribed rows \textup{(}\resp
columns\textup{)} in the chosen block row \textup{(}\resp
column\textup{)}. Then there exists an $\bA$ in $\E_\Bbbk(n,r)$
satisfying $\rho(\bA)=\bB$ which agrees with the prescribed rows
\textup{(}\resp columns\textup{)}.
\end{lemm}

\begin{proof}
Suppose we have fixed on an $i$th block row extension (the argument
for a block column extension is similar and left to the reader). By
hypothesis, a set $D(n,r-1)$ exists satisfying Property~\ref{4} for
$(n,r-1)$. Any assignment to this set determines a decomposition of
the given $\bB$, say $\bB = \bB(1)+ \cdots + \bB(n)$ where each
$\bB(j) \in \E_\Bbbk(n,r-1)^i_j$. We define $F'(n,r)$ to correspond
to $D(n,r-1)$ under the bijection $\pi^i$ of Lemma~\ref{lem:bijections}. Assignments to $D(n,r-1)$ correspond to
assignments to $F'(n,r)$, which complete the $i$th block row of $\bA$
in a way that satisfies all relevant extension equations. Under this
correspondence, $\bA^i_j = \bB(j)$ for each $j=1,\dots,n$. Once the
$i$th block row of $\bA$ is complete, we apply Proposition~\ref{prop:extend}\ref{prop6.1_b} to find special extensions $\bA(j)$ in $\E_\Bbbk(n,r)^i_j$ for all $j$ such that $\bA=\bA(1)+ \cdots +
\bA(n)$ extends $\bB$.

To finish, we simply note that by hypothesis the prescribed rows in
the $i$th block row are compatible with the extension problem for the
given $\bB$, so they are specified by a partial assignment $f'$ to
some subset of some $F'(n,r)$. We can extend $f'$ to an assignment
$f:F'(n,r)\to \Bbbk$ which then determines the $i$th block row of
$\bA$ as above, in a way that coincides with the prescribed
information.
\end{proof}



\goodbreak
\begin{exam}
We illustrate the above proof. Take $(n,r)=(4,2)$, and suppose that
row 42 of an extension has been prescribed, for a given $\bB$ in
$\E_\Bbbk(4,1)$. The prescribed row is part of the final block row, so
we construct a final block row extension. The following
\[
F'(4,2): \quad
\text{\setlength{\arraycolsep}{2.7pt}\scriptsize
$\displaystyle \begin{array}{c|ccc|ccc|ccc|ccc|}
 &\rt{12}&\rt{13}&\rt{14}&\rt{21} &\rt{23}&\rt{24}&\rt{31}&\rt{32}
 &\rt{34}&\rt{41}&\rt{42}&\rt{43} \\ \hline 
 42 &&&		&&&\cm		&&\cm&\cm	&&\cm&\cm\\ 
 43 &&&		&&&\cm		&&\cm&\cm	&&\cm&\cm\\ 
 \hline
\end{array}\; $}
\]
depicts a row- and column-terminal free-pattern $F'(4,2)$ for the last
block row of our desired extension $\bA$, where the checkmarked entries
correspond to elements of $F'(4,2)$, which can be freely assigned in
order to determine the last block row of a general extension. This
free pattern is calculated in Example~\ref{Dn1} below.

The prescribed $42$-row merely determines the values of the five free
entries in that row. By assigning arbitrary values to the remaining
five entries in $F'(4,2)$ we complete the last block row of $\bA$, and
then complete $\bA$ by choosing special extensions of each block in
the last block row, and summing, as in the first paragraph of the
proof of Lemma~\ref{lem:partial}.
\end{exam}


We note the following immediate consequence of Lemma~\ref{lem:partial}, which was promised in Remark~\ref{rmk:promise}.

\begin{coro}\label{cor:colouring}
Assume the same hypotheses as in Lemma~\ref{lem:partial}. Suppose that
$(a^\bi_*)$ or $(a^*_\bj)$ is a row or column \textup{(}where $\bi$ or
$\bj \in I'(n,r)$\textup{)} determined by an assignment to the
variables in that row or column labelled by the set $I_1(n,r)$ in
Algorithm~\ref{algo}. Then that row or column is a row or column of
some extension $\bA$ of any given $\bB$ in $\E_\Bbbk(n,r-1)$.
\end{coro}

Now we are ready to prove Property~\ref{2}. 

\begin{prop}\label{prop:decomp}
Assume Property~\ref{1} for $(n-1,r)$ and Property~\ref{4} for
$(n-1,r-1)$. Then Property~\ref{2} holds for $(n,r)$.
\end{prop}


\begin{proof}
Let $\bA$ in $\E_\Bbbk(n,r)$ be given. We choose to decompose $\bA$
based on its last block row. (The argument for any other block row or
column is similar.) By Proposition~\ref{prop:n-reduce-iso}, the
problem of extending from $\E_\Bbbk(n,r-1)^n_j$ to $\E_\Bbbk(n,r)^n_j$
is equivalent to the problem of extending from $\E_\Bbbk(n-1,r-1)$ to
$\E_\Bbbk(n-1,r)$. So by the first hypothesis and Proposition~\ref{prop:extend}\ref{prop6.1_a},
\begin{equation}\label{eq:spec-ext}
\text{there exists $\bA(j)$ in $\E_\Bbbk(n,r)^n_j$ such that
 $\rho(\bA(j)) = \bA^n_j$}
\end{equation}
for all $j = 1, \dots, n$. Any choice of $\bA(j)$ satisfying~\eqref{eq:spec-ext} makes the last block row of $\bA$ agree with that of the sum $\bA(1)+\cdots +\bA(n)$, so that the last block row of the
difference $\bC=\bA-\bA(1)-\cdots -\bA(n)$ is zero. We need to show
that it is always possible to choose the $\bA(j)$ satisfying~\eqref{eq:spec-ext} in such a way that $\bC=\bzero$.

\begin{enonce*}{\hypertarget{Case 1}{Case 1}} If $n \le r+1$ then $\bC$ as
above is an invariant whose last block row is zero. The zero invariant
$\bzero$ is another such invariant. By Remark~\ref{rmk:sharper}, it
follows that $\bC= \bzero$. This completes \hyperlink{Case 1}{Case 1}.
\end{enonce*}

\begin{enonce*}{\hypertarget{Case 2}{Case 2}} Assume for the rest of the proof that $n > r+1$ (so
$n-r \ge 2$). We aim to show that $\bA(j)$ for $j=2, \dots, n$ can be
chosen satisfying~\eqref{eq:spec-ext} in such a way that all but the
last block of the first block column of $\bA - \bA(2) - \cdots -
\bA(n)$ is zero. Working in \emph{reverse} order from right to left
along the last block row of $\bA$, we choose $\bA(j)$ for $j = n, n-1,
\dots, 2$ satisfying~\eqref{eq:spec-ext} by the following process. For
each $j$, assuming that $\bA(n), \dots, \bA(j)$ have already been
chosen, we set $\bC(j)=\bA-\bA(n)-\cdots -\bA(j)$.

\begin{enonce*}[remark]{\hypertarget{Step 1}{Step 1}} First we choose arbitrary $\bA(n), \dots, \bA(r+2)$
subject only to condition~\eqref{eq:spec-ext}. This step is not
inductive.
\end{enonce*}

\begin{enonce*}[remark]{\hypertarget{Step 2}{Step 2}} We proceed by reverse induction on $j$ running from
$r+1$ down to $2$. At each stage, we claim that $\bA(j)$ satisfying~\eqref{eq:spec-ext} can be chosen so that:
\begin{equation}\label{eq:agrees}
 \begin{minipage}{4in}
 $\eta^n_j(\bA(j))$ agrees with $\eta^n_j(\bC(j+1))$ on all columns
 indexed by a label in $\Lin_{j-1}$
 \end{minipage}
\end{equation}
where $\Lin_j = \{j_1\cdots j_r \in I'(n,r): j_\alpha=\alpha \text{
 for all } \alpha=1, \dots, j\}$. To see this, we apply Lemma~\ref{lem:partial} in the case $(n-1,r)$, since $\eta^n_j(\bC(j+1))$
belongs to $\E_\Bbbk(n-1,r)$. (Property~\ref{1} for $(n-1,r)$ implies
the existence of special extensions for $(n-1,r)$, which is equivalent
to Property~\ref{1} for $(n-2,r)$, so the hypotheses of Lemma~\ref{lem:partial} for the case $(n-1,r)$ are satisfied.) By the lemma, we can find $\bA'(j)$ in $\E_\Bbbk(n-1,r)$ which agrees on the image
under $\eta^n_j$ of the columns of $\bC(j)$ indexed by
$\Lin_{j-1}$. Then we set $\bA(j) = \theta^n_j \bA'(j)$. This matrix
satisfies~\eqref{eq:agrees}, so the claim is proved.

At this point we have inductively chosen $\bA(r+1)$ down to $\bA(2)$,
and we are ready for the final step, choosing $\bA(1)$. Since
$\Lin_1$ indexes all the non-initialised columns in the first block
column, the only nonzero block in the first block column of $\bC(2)$
is $\bC(2)^n_1$. By construction, the same is true of the blocks in
the last block row. Thus, we may apply Lemma~\ref{lem:special-criterion} to conclude that $\bC(2)$ belongs to $\E_\Bbbk(n,r)^n_1$, and hence by setting $\bA(1) = \bC(2)$ we are
done.\qedhere
\end{enonce*}
\end{enonce*}
\end{proof}


\hyperlink{Step 2}{Step 2} above starts with the unique column in $\Lin_r = \{1 \cdots
r\}$, which isn't affected by any special invariant in
$\E_\Bbbk(n,r)^n_j$ for $j<r+1$. So the last opportunity to zero that
column is when we choose $\bA(r+1)$. Similarly, as the induction in
\hyperlink{Step 2}{Step~2} proceeds, controlled by the nested sequence
\[
\Lin_{r} \subset \Lin_{r-1} \subset \cdots \subset \Lin_1,
\]
the new columns that are zeroed in the running difference are
precisely those columns that cannot be affected in subsequent steps. 

 

\begin{exam}\label{ex:decomp42}
Assume that $\bA \in \E_\Bbbk(4,2)$. We illustrate \hyperlink{Case 2}{Case~2} of the above
proof for a last block row decomposition, working from right to left
through the last block row. \hyperlink{Step 1}{Step~1} consists of subtracting an
arbitrary $\bA(4)$ in $\E_\Bbbk(4,2)^4_4$ such that
$\rho(\bA(4))=\bA^4_4$. Referring to the matrix forms depicted in
Examples~\ref{ex:n23-r2-inv},~\ref{ex:special-inv} we see that after
\hyperlink{Step 1}{Step~1},
\[
\bC(4) = \bA - \bA(4) = 
\text{\setlength{\arraycolsep}{2.7pt}\scriptsize
$\displaystyle 
\begin{array}{c|cccc|cccc|cccc|cccc|}
 &\rt{11}&\rt{12}&\rt{13}&\rt{14}&\rt{21}&\rt{22}&\rt{23}&\rt{24}&
 \rt{31}&\rt{32}&\rt{33}&\rt{34}&\rt{41}&\rt{42}&\rt{43}&\rt{44}\\ \hline
 11&*& & & & &*& & & & &*& & & & &*\\
 12& &t_1&*&*&*& &*&*&*&*& &*&*&*&* & \\
 13& &t_2&*&*&*& &*&*&*&*& &*&*&*&*& \\
 14& &*&s&0&*& &*&0&*&*& &0&*&*&*& \\ \hline
 21& &t'_1&*&*&*& &*&*&*&*& &*&*&*&*& \\ 
 22&*& & & & &*& & & & &*& & & & &*\\
 23& &t'_2&*&*&*& &*&*&*&*& &*&*&*&*& \\ 
 24& &*&s'&0&*& &*&0&*&*& &0&*&*&*& \\ \hline 
 31& &t''_1&*&*&*& &*&*&*&*& &*&*&*&*& \\
 32& &t''_2&*&*&*& &*&*&*&*& &*&*&*&*& \\
 33&*& & & & &*& & & & &*& & & & &*\\ 
 34& &*&s''&0&*& &*&0&*&*& &0&*&*&*& \\ \hline
\end{array} $}
\]
where blank entries represent zero as usual. We omit showing the last
block row, where there are no choices. The nine explicit zeros shown
above are place-permutation symmetric to entries of $\bC(4)^4_4$,
which is $\bzero$ by construction. We must have $t_1+t_2 = s$,
$t'_1+t'_2 = s'$, and $t''_1+t''_2 = s''$ thanks to local GDS
conditions in the blocks.

The above equations imply that the image of column 12 under $\eta^4_3$
is consistent with that column of an extension of
$\eta^4_3(\bA^4_3)$. By Lemma~\ref{lem:partial}, we can find an
$\bA(3)$ in $\E_\Bbbk(4,2)^4_3$ such that $\rho(\bA(3)) = \bC(4)^4_3$
and $\eta^4_3 (\bA(3))$ agrees with $\eta^4_3(\bC(4))$ in column~12. 
(Note that $\Lin_2=\{12\}$.) This implies that $\bC(3)$ has the form
\[
\bC(3) = \bC(4) - \bA(3) = 
\text{\setlength{\arraycolsep}{2.7pt}\scriptsize
$\displaystyle 
\begin{array}{c|cccc|cccc|cccc|cccc|}
 &\rt{11}&\rt{12}&\rt{13}&\rt{14}&\rt{21}&\rt{22}&\rt{23}&\rt{24}&
 \rt{31}&\rt{32}&\rt{33}&\rt{34}&\rt{41}&\rt{42}&\rt{43}&\rt{44}\\ \hline
 11&*& & & & &*& & & & &*& & & & &*\\
 12& &0&*&*&0& &*&*&*&*& &*&*&*&* & \\
 13& &0&*&*&0& &*&*&*&*& &*&*&*&*& \\
 14& &*&0&0&*& &0&0&*&*& &0&*&*&0& \\ \hline
 21& &0&*&*&0& &*&*&*&*& &*&*&*&*& \\ 
 22&*& & & & &*& & & & &*& & & & &*\\
 23& &0&*&*&0& &*&*&*&*& &*&*&*&*& \\ 
 24& &*&0&0&*& &0&0&*&*& &0&*&*&0& \\ \hline 
 31& &0&*&*&0& &*&*&*&*& &*&*&*&*& \\
 32& &0&*&*&0& &*&*&*&*& &*&*&*&*& \\
 33&*& & & & &*& & & & &*& & & & &*\\ 
 34& &*&0&0&*& &0&0&*&*& &0&*&*&0& \\ \hline
% 41& &* &* &* &*& &*&*&&& &&&&& \\
% 42& &* &* &* &*& &*&*&&& &&&&& \\
% 43& &* &* &* &*& &*&*&&& &&&&& \\
% 44&* & & & & &*& & & & && & & & &\\ \hline
 \end{array} \; .$}
\]
The explicit zeros in column 21 must be zero because invariants are
place-permutation invariant, and column 21 is place-permutation
symmetric to column 12.

By Lemma~\ref{lem:partial} there exists an $\bA(2)$ in
$\E_\Bbbk(4,2)^4_2$ such that $\rho(\bA(2)) = \bA^4_2$ and $\eta^4_2
(\bA(2))$ agrees with $\eta^4_2(\bC(3))$ in columns $\Lin_1 = \{1*\} =
\{12, 13, 14\}$. Hence the difference $\bC(2) = \bC(3) - \bA(2)$
satisfies the property
\[
\text{$\bC(2)^i_4 = \bzero$ and $\bC(2)^4_j = \bzero$ for all $i, j
 <4$.}
\]
By Lemma~\ref{lem:special-criterion} it follows that $\bC(2)$ is
in $\E_\Bbbk(4,2)^4_1$, so by setting $\bA(1) = \bC(2)$ we obtain the
desired decomposition $\bA = \bA(1)+ \cdots + \bA(4)$. This completes
Example~\ref{ex:decomp42}.
\end{exam}

To prove Property~\ref{3} we need the following result, which
describes how to construct a free pattern for the special extension
problem, for a given pair $i,j$ of indices in the set $\{1, \dots,
n\}$. We remind the reader that, by Proposition~\ref{prop:special-ext-form}\ref{prop4.8_c}, the restriction of any $\bA$ in
$\E_\Bbbk(n,r)^i_j$ belongs to $\E_\Bbbk(n,r-1)^i_j$. Let $\theta^i_j$
be the isomorphism in Proposition~\ref{prop:n-reduce-iso}.

\begin{lemm}\label{lem:free-special}
Assume Property~\ref{3} for $(n-1,r)$. Then $F(n,r)^i_j =
\theta^i_jF(n-1,r)$ is a free pattern for the problem of extending
invariants from $\E_\Bbbk(n,r-1)^i_j$ to $\E_\Bbbk(n,r)^i_j$, in the
sense that for any given $\bB$ in $\E_\Bbbk(n,r-1)^i_j$ and any
assignment $f^i_j: F(n,r)^i_j \to \Bbbk$, there is a unique special
invariant $\bA$ in $\E_\Bbbk(n,r)^i_j$ such that $\rho(\bA) = \bB$.
\end{lemm}

\begin{proof}
Suppose some $\bB$ in $\E_\Bbbk(n,r-1)^i_j$ is given. To construct a
special extension $\bA$ in $\E_\Bbbk(n,r)^i_j$ we must set $\bA^i_j =
\bB$ and all other blocks in the $i$th block row and $j$th block
column to $\bzero$. The rest of $\bA$ is determined by
place-permutation symmetry and the isomorphism $\theta^i_j:
\E_\Bbbk(n-1,r) \to \E_\Bbbk(n,r)^i_j$ of Proposition~\ref{prop:n-reduce-iso}, thus determined by assigning values to images
of the free variables in the free pattern $F(n-1,r)$.
\end{proof}


Now we are ready for the proof of Property~\ref{3}.

\begin{prop}\label{prop:free-ptns-exist}
 Property~\ref{4} for $(n,r-1)$ and Property~\ref{3} for $(n-1,r)$
 imply Property~\ref{3} for $(n,r)$.
\end{prop}


\begin{proof} 
We choose to base the construction of $F(n,r)$ on the last block
row. The argument for any other block row or column is similar. Let
$\pi = \pi^n$ be the map in Lemma~\ref{lem:bijections}. Set $F'(n,r)
= \pi^{-1} D(n,r-1)$. Then $F'(n,r)$ is a free pattern for completing
the last block row of an extension $\bA$ of $\bB$, because doing so is
equivalent to decomposing $\bB$ along its last block row. Let $\{
\bA^n_j: j=1, \dots, n \}$ be the last block row determined by some
chosen assignment $f': F'(n,r) \to \Bbbk$. Fix this choice of last
block row of $\bA$ for the rest of the argument.

To find the desired extension $\bA$, we apply Proposition~\ref{prop:extend}\ref{prop6.1_a} to choose arbitrary special extensions $\bA(j)$
in $\E_\Bbbk(n,r)^n_j$ such that $\rho(\bA(j)) = \bA^n_j$ for all $j$,
and set $\bA = \bA(1) + \cdots + \bA(n)$. By Lemma~\ref{lem:free-special}, there exists a free pattern $F(n,r)^n_j$ for
each $j$ and an assignment $f(j): F(n,r)^n_j \to \Bbbk$ that
determines $\bA^n_j$. Let $F''(n,r) = \bigcup_{j=1}^n F(n,r)^n_j$. To
continue, we introduce the notation
\[
\bA_{F''} := (a^\bi_\bj)_{(\bi,\bj) \in F''(n,r)}
\]
for the restriction of $\bA$ to $F''(n,r)$, and similarly for each
$\bA(j)$. We have $\bA_{F''} = \sum_j \bA(j)_{F''}$, so the given
assignments $f(j)$ induce a corresponding assignment $f'': F''(n,r)
\to \Bbbk$. For the given fixed last block row of $\bA$, it is clear
that there is an extension $\bA$ whose restriction to $F''(n,r)$
induces $f''$, for an arbitrary assignment $f'': F''(n,r) \to \Bbbk$,
because we can choose the $f(j)$ entries arbitrarily.

Every extension $\bA$ of $\bB$ with the specified last block row must
be of the form $\bA = \bA(1)+\cdots +\bA(n)$, where $\bA(j)$ is in
$\E_\Bbbk(n,r)^n_j$ and $\rho(\bA(j)) = \bA^n_j$ for all $j$. This
follows from the hypotheses, which by Proposition~\ref{prop:decomp}
imply that the decomposition property holds for $(n,r)$.

To finish, we need to argue that $\bA$ (with the fixed last block row)
is uniquely determined by its assignment $f''$. To see this, suppose
that $\bA'$ is another extension of $\bB$, with the same last block
row, given by the same assignment $f''$. Then $\bA_{F''} =
\bA'_{F''}$. Thus the equations $\bA = \sum_j \bA(j)$, $\bA' = \sum_j
\bA'(j)$ imply by restriction that
\[
\textstyle \sum _j \bA(j)_{F''} = \sum_j \bA'(j)_{F''} \, .
\]
Each entry of $\bA(j)$, $\bA'(j)$ is uniquely expressible by the same
linear combination of the free pattern variables $a(j)^\bi_\bj$,
$a'(j)^\bi_\bj$ indexed by the free pattern $F(n,r)^n_j$. Hence, any
linear relations among the $\{\bA(j)\}$, $\{\bA'(j)\}$ are determined
by their restriction to $F''(n,r)$, which contains $F(n,r)^n_j$. So
the above displayed equality implies that
\[
\textstyle \sum _j \bA(j) = \sum_j \bA'(j) \, .
\]
Hence $\bA = \bA'$, as required. This proves the desired uniqueness
statement. We have now shown that $F''(n,r)$ is a free pattern for
constructing an extension $\bA$ having the specified last block
row. Thus, the disjoint union $F(n,r) = F'(n,r) \sqcup F''(n,r)$ is a
free pattern for the extension problem (for the given $\bB$).
\end{proof}

We note the following immediate consequence of the above proof.

\begin{coro}
 Under the same hypotheses as the preceeding result based on an $i$th
 block row \textup{(}respectively, $j$th block column\textup{)}
 construction,
 \[
 F(n,r) = F'(n,r) \sqcup F''(n,r)
 \]
 where $F''(n,r) = \bigcup_{j=1}^n F(n,r)^i_j$ \textup{(}\resp
 $F''(n,r) = \bigcup_{i=1}^n F(n,r)^i_j$\textup{)} and $F'(n,r)$ is
 determined by the condition $\pi F'(n,r) = D(n,r-1)$.
\end{coro}


\begin{exam}\label{ex:free-patterns}\ \\*[-1.2em]
\begin{enumerate}[label = (\alph*)]
\item \label{exam6.12_i} 
We now construct $F(4,2)$, under the assumption that $F(3,2)$ and
 $D(4,1) \cong F'(4,2)$ are known. It is easy to check that $F(3,2)
 = \{(32,32)\}$. Hence
\[
F''(4,2)=\bigcup_{1\leq j \leq 4}\theta^j_4(F(3,2)) =\bigcup_{1\leq j
 \leq 4}\theta^j_4\{(32),(32)\}
\]
where $\theta^1_4\{32,32\}= \{32,43\}$, $\theta^2_4\{32,32\}=
\{32,43\}$, $\theta^3_4\{32,32\}= \{32,42\}$, and
$\theta^4_4\{32,32\}= \{32,32\}$. We refer to Example~\ref{Dn1} for
$F'(4,2)$. It follows that $F(4,2)$ is as depicted below:
 \[ F(4,2): \quad
\text{\setlength{\arraycolsep}{2.7pt}\scriptsize
$\displaystyle 
\begin{array}{c|ccc|ccc|ccc|ccc|}
 &\rt{12}&\rt{13}&\rt{14}&\rt{21} &\rt{23}&\rt{24}&\rt{31}&\rt{32}
 &\rt{34}&\rt{41}&\rt{42}&\rt{43} \\ \hline
 32 &&&	&&&	&&\cm&	&&\cm&\cm\\ \hline
 42 &&&	&&&\cm	&&\cm&\cm	&&\cm&\cm\\ 
 43 &&&	&&&\cm	&&\cm&\cm	&&\cm&\cm\\ 
 \hline
\end{array}\; .$}
\]

\item \label{exam6.12_ii} 
Now we consider $F(5,3)$, under the assumption that $F(4,3)$ and
$D(5,2)$ are both known. It is easy to check that $F(4,3) = \{(432,
432)\}$. It follows that
 \[
 F'' (5,3): \quad
\text{\setlength{\arraycolsep}{2.7pt}\scriptsize
$\displaystyle 
\begin{array}{c|c|ccc|ccccc|ccccc|}
 &\rt{254}&\rt{325}&\rt{352}&\rt{354}%
 &\rt{425}&\rt{432}&\rt{435}&\rt{452}%
 &\rt{453}&\rt{524}&\rt{532}&\rt{534}%
 &\rt{542}&\rt{543}\\ \hline
 432 & && & &&\cm&& & &&\cm&&\cm&\cm \\ 
 \hline
\end{array}\,$}
 \]
since $\theta^5_5(432,432)=(432,432)$,
$\theta^4_5(432,432)=(432,532)$, $\theta^3_5(432,432)=(432,542)$, and
$\theta^2_5(432,432)=(432,543)$. See Example~\ref{ex:D52} for
$F'(5,3)$. Taking the union of $F''(5,3)$ with $F'(5,3)$ gives
$F(5.3)$ as depicted below:
\[
 F(5,3): \quad
\text{\setlength{\arraycolsep}{2.7pt}\scriptsize
$\displaystyle 
\begin{array}{c|c|ccc|ccccc|ccccc|}
 &\rt{254}&\rt{325}&\rt{352}&\rt{354}%
 &\rt{425}&\rt{432}&\rt{435}&\rt{452}%
 &\rt{453}&\rt{524}&\rt{532}&\rt{534}%
 &\rt{542}&\rt{543}\\ \hline
 432 & && & &&\cm&& & &&\cm&&\cm&\cm \\ \hline
 532 &\cm &&\cm&\cm &&\cm&&\cm&\cm &&\cm&&\cm&\cm \\ 
 542 &\cm &\cm&\cm&\cm &\cm&\cm&\cm&\cm&\cm &\cm&\cm&\cm&\cm&\cm \\
 543 &\cm &\cm&\cm&\cm &\cm&\cm&\cm&\cm&\cm &\cm&\cm&\cm&\cm&\cm \\ \hline
\end{array}\,.$}
\]
This ends Example~\ref{ex:free-patterns}.
\end{enumerate}
\end{exam}




It remains to prove Property~\ref{4}. The proof given below closely
follows the proof of Proposition~\ref{prop:decomp}.


\begin{prop}\label{prop:property-4}
Property~\ref{3} for $(n-1,r)$ implies Property~\ref{4} for $(n,r)$.
\end{prop}

\begin{proof}
We need to prove the existence of the set $D(n,r)$ satisfying
Property~\ref{4}. Property~\ref{3} for $(n-1,r)$ implies
Property~\ref{4} for $(n-1,r-1)$ and also implies Property\ref{1} for
$(n-1,r)$, so the hypotheses of Proposition~\ref{prop:decomp} are
satisfied, and hence its conclusion holds. In other words, any given
$\bA$ in $\E_\Bbbk(n,r)$ has a decomposition based on a chosen block
row or column. We assume for concreteness that it is based on the last
block row, as in the proof of Proposition~\ref{prop:decomp}. The
argument closely follows the proof of that proposition.

\begin{enonce*}[remark]{\hypertarget{prop6.13_C1}{Case 1}} If $n \le r+1$ then by Case 1 of
the proof of Proposition~\ref{prop:decomp}, the desired decomposition
$\bA = \bA(1) + \cdots+ \bA(n)$ is unique. Hence $D(n,r)$ is empty in
this case.
\end{enonce*}

\begin{enonce*}[remark]{\hypertarget{prop6.13_C1}{Case 2}} Assume henceforth that $n > r+1$.
By Lemma~\ref{lem:free-special} and the hypothesis, free patterns
\[F(n,r)^n_j = \theta^n_j F(n-1,r)\] are available for any
$j=1,\dots, n$. We may assume that $F(n-1,r)$ is row- and
column-terminal.
\end{enonce*}

\begin{enonce*}[remark]{\hypertarget{prop6.13_S1}{Step 1}} For each $j=r+2, \dots, n$ we choose arbitrary
assignments $F(n,r)^n_j \to \Bbbk$. Each assignment uniquely
determines a matrix $\bA(j)$ in $\E_\Bbbk(n,r)_j^n$. Therefore we
define
\begin{equation*}
D'(n,r) = \bigcup_{r+2 \le j \le n} \{j\}\times F(n,r)^n_j \, .
\end{equation*}
This set parametrises a set of free entries that uniquely determines
matrices $\bA(r+2), \dots, \bA(n)$ in~\hyperlink{Case 2}{Case 2} of the proof of
Proposition~\ref{prop:decomp}. Set $\bC(j+2) = \bA - \sum_{j=r+2}^n
\bA(j)$.
\end{enonce*}

\begin{enonce*}[remark]{\hypertarget{prop6.13_S1}{Step 2}} Now we proceed by reverse induction on $j$
running from $r+1$ down to $2$, in order to identify a set $D''(n,r)$
of free entries parametrising the choice of $\bA(j+1), \dots, \bA(2)$
satisfying condition~\eqref{eq:agrees} in the proof of Proposition~\ref{prop:decomp}. Clearly the union
\begin{equation*}\label{toomuch2}
\bigcup_{2\leq j \leq r+1} \{j\}\times F(n,r)^n_j 
\end{equation*}
is an upper bound on $D''(n,r)$. This bound is not tight due to linear
dependencies caused by the prescription of columns labelled by
elements of ${\Lin_{j-1}}^{\Sym_r}$ in the proof of~\hyperlink{Step 2}{Step 2} of
Proposition~\ref{prop:decomp}. We have to remove those additional
dependencies in order to obtain $D''(n,r)$.

We do this by applying a modified version of Algorithm~\ref{algo}. Namely, for each fixed $j = r+1, \dots, 2$ we initialise
each entry $j_1 \cdots j_r$ of $I'(n,r)$ containing $j$ to colour 0,
and (following ~\hyperlink{Step 2}{Step 2} in~\hyperlink{Case 2}{Case 2} of the proof of Proposition~\ref{prop:decomp}) do the same for each entry belonging to
${\Lin_{j-1}}^{\Sym_r}$. Then we run Algorithm~\ref{algo} to determine
the independent (free) columns in a generic row of an extension. Let
$I'_j(n,r)$ be the set of elements coloured 1 in that algorithm. We
define
\[
\ov{F}(n,r)^n_j = \{(i_1\cdots i_r, j_1\cdots j_r) \in F(n,r)^n_j :
j_1\cdots j_r \in I'_j(n,r)\}.
\]
and we accordingly set
\[
D''(n,r) = \bigcup_{2\leq j \leq r+1} \{j\}\times
 \ov{F}(n,r)^n_j \, .
\]
We claim that the desired pattern $D(n,r)$ is equal to the union
\[
D(n,r) = D'(n,r) \cup D''(n,r).
\]
To see this, observe that every decomposition $\bA = \bA(1) +\cdots+
\bA(n)$ in Proposition~\ref{prop:decomp} determines a unique
assignment
\begin{equation}\label{eq:assign-D}
f: D(n,r) \to \Bbbk
\end{equation}
by setting $f(\{j\} \times (\bp,\bq)) = a(j)^\bp_\bq$. Conversely,
for each $r+1 \ge j \ge 2$, each assignment to $\ov{F}(n,r)^n_j$ along
with the values in the prescribed columns indexed by
${\Lin_{j-1}}^{\Sym_r}$ forces a corresponding assignment to
$F(n,r)^n_j$. (Indeed, the algorithm was designed with that purpose in
mind.) It thus follows from the proof of~\hyperlink{Case 2}{Case 2} of Proposition~\ref{prop:decomp} that each assignment as in~\eqref{eq:assign-D}
determines a unique decomposition of the form $\bA = \bA(1) +\cdots+
\bA(n)$ with the required properties.\qedhere
\end{enonce*}
\end{proof}

\begin{exam}\label{Dn1}
To illustrate the above proof, we construct $D(n,1)$, assuming that
$F(n-1,1) = \{(i,j): 2 \le i,j \le n-1\}$. For each $j=1, \dots,
n$ we have
\[
F(n,1)^n_j = \{ (p,q): 2 \le p \le n-1,\ 2 \le q \le n,\ q \ne j\}.
\]
Furthermore, $\ov{F}(n,1)^n_2$ is obtained from $F(n,1)^n_2$ by
excising all entries in its leftmost column; that is,
\[
\ov{F}(n,1)^n_2 = \{ (p,q): 2 \le p \le n-1,\ 4 \le q \le n\}.
\]
Thus, as in the proof of Proposition~\ref{prop:property-4}, we obtain
\[
D(n,1) = \{2\}\times \ov{F}(n,r)^n_2 \quad\cup\quad \bigcup_{3\leq j
 \leq n}\{j\}\times F(n,r)^n_j \, .
\]
For instance, when $n=4$ this can be depicted by the following table:
\[ D(4,1): \quad
\text{\setlength{\arraycolsep}{2.7pt}\scriptsize
$\displaystyle 
\begin{array}{c|cccc||cccc||cccc||cccc|}
\hline
j=&\multicolumn{4}{c||}{1}	
&\multicolumn{4}{c||}{2}
&\multicolumn{4}{c||}{3}
&\multicolumn{4}{c|}{4}					\\ \hline
 &1&2&3&4 &1&2&3&4 &1&2&3&4 &1&2&3&4 \\ \hline
 1 &&&& &&&& &&&& &&&&\\
 2 &&&& &&&&\cm && \cm&&\cm && \cm&\cm&\\
 3 &&&& &&&&\cm && \cm&&\cm && \cm&\cm&\\ 
 4 &&&& &&&& &&&& &&&&\\ 
 \hline
\end{array}$}
\]
in which the indices labelling checkmarked positions correspond to
elements of $D(4,1)$. By Lemma~\ref{lem:bijections}, it follows that
$F'(4,2)$ can be depicted by
\[
F'(4,2): \quad
\text{\setlength{\arraycolsep}{2.7pt}\scriptsize
$\displaystyle 
\begin{array}{c|ccc|ccc|ccc|ccc|}
 &\rt{12}&\rt{13}&\rt{14}&\rt{21} &\rt{23}&\rt{24}&\rt{31}&\rt{32}
 &\rt{34}&\rt{41}&\rt{42}&\rt{43} \\ \hline
% 32 &&&	&&&	&&\cm&	&&\cm&\cm\\ \hline
 42 &&&	&&&\cm	&&\cm&\cm	&&\cm&\cm\\ 
 43 &&&	&&&\cm	&&\cm&\cm	&&\cm&\cm\\ 
 \hline
\end{array}\; .$}
\]
\end{exam}


\begin{exam}\label{ex:D52}
We now compute $D(5,2)$, which is in bijection with $F'(5,3)$. We
first run Algorithm~\ref{algo} to calculate
$I'_j(5,2)$ for $j=3,2$:
\[
 I'_3(5,2): \quad
\text{\setlength{\arraycolsep}{2.7pt}\scriptsize
$\displaystyle 
\begin{array}{c|cccc|cccc|cccc|cccc|cccc|}
 &\rt{12}&\rt{13}&\rt{14}&\rt{15}&\rt{21} &\rt{23}&\rt{24}&\rt{25}
 &\rt{31}&\rt{32} &\rt{34}&\rt{35}&\rt{41}&\rt{42}&\rt{43} &\rt{45} 
 &\rt{51}&\rt{52}&\rt{53} &\rt{54} \\ \hline
 &0&0&0	&0	&0&0&0	&1	 &0&0&0&0	&0&1&0&1 &0&1 &0&1
 \\ 
 \hline
\end{array}\; $} 
\]
\[
 I'_2(5,2): \quad
\text{\setlength{\arraycolsep}{2.7pt}\scriptsize
$\displaystyle 
\begin{array}{c|cccc|cccc|cccc|cccc|cccc|}
 &\rt{12}&\rt{13}&\rt{14}&\rt{15}&\rt{21} &\rt{23}&\rt{24}&\rt{25}&\rt{31}&\rt{32} &\rt{34}&\rt{35}&\rt{41}&\rt{42}&\rt{43} &\rt{45} 
 &\rt{51}&\rt{52}&\rt{53} &\rt{54} \\ \hline
 &0&0&0	&0	 &0&0&0&0&0&0&0	&1	&0&0&1&1	&0&0 &1&1
 \\ 
 \hline
\end{array}\;.$} 
\]
To justify this, we go through the algorithm for $I'_3(5,2)$ a little
more slowly. The initial colouring is
 \[ 
\text{\setlength{\arraycolsep}{2.7pt}\scriptsize
$\displaystyle 
\begin{array}{c|cccc|cccc|cccc|cccc|cccc|}
 &\rt{12}&\rt{13}&\rt{14}&\rt{15}&\rt{21} &\rt{23}&\rt{24}
 &\rt{25}&\rt{31}&\rt{32} &\rt{34}&\rt{35}&\rt{41}&\rt{42}&\rt{43} &\rt{45} 
 &\rt{51}&\rt{52}&\rt{53} &\rt{54} \\ \hline 
 &0&0&	&	&0&0&&	&0	&0&0&0	&	&&0&	& &&0&	\\ 
 \hline
\end{array}\; .$}
\] 
We now list the colourings obtained after three complete iterations
(doing both~\hyperlink{prop6.13_S1}{Steps~1} and~\hyperlink{prop6.13_S2}{2}) of the algorithm:
 \[ 
\text{\setlength{\arraycolsep}{2.7pt}\scriptsize
$\displaystyle 
\begin{array}{c|cccc|cccc|cccc|cccc|cccc|}
 &\rt{12}&\rt{13}&\rt{14}&\rt{15}&\rt{21} &\rt{23}&\rt{24}&\rt{25}&\rt{31}&\rt{32} &\rt{34}&\rt{35}&\rt{41}&\rt{42}&\rt{43} &\rt{45} 
 &\rt{51}&\rt{52}&\rt{53} &\rt{54} \\ \hline 
 &0&0&		&	&0&0&&	&0	&0&0&0	&	&&0&	&		&&0&1			\\ 
 \hline
\end{array}\;$} 
\] 
 \vspace{1mm}
 \[ 
\text{\setlength{\arraycolsep}{2.7pt}\scriptsize
$\displaystyle 
\begin{array}{c|cccc|cccc|cccc|cccc|cccc|}
 &\rt{12}&\rt{13}&\rt{14}&\rt{15}&\rt{21} &\rt{23}&\rt{24}&\rt{25}&\rt{31}&\rt{32} &\rt{34}&\rt{35}&\rt{41}&\rt{42}&\rt{43} &\rt{45} 
 &\rt{51}&\rt{52}&\rt{53} &\rt{54} \\ \hline 
 &0&0&		&	&0&0&&	&0	&0&0&0	&0	&0&0&0	&0&1&0&1			\\ 
 \hline
\end{array}\; $} 
\] \vspace{1mm}
 \[ 
\text{\setlength{\arraycolsep}{2.7pt}\scriptsize
$\displaystyle 
\begin{array}{c|cccc|cccc|cccc|cccc|cccc|}
 &\rt{12}&\rt{13}&\rt{14}&\rt{15}&\rt{21} &\rt{23}&\rt{24}&\rt{25}&\rt{31}&\rt{32} &\rt{34}&\rt{35}&\rt{41}&\rt{42}&\rt{43} &\rt{45} 
 &\rt{51}&\rt{52}&\rt{53} &\rt{54} \\ \hline 
 &0&0&0		&0	&0&0&0&1	&0	&0&0&0	&0	&0&0&0	&0&1&0&1			\\ 
 \hline
\end{array}\;$} 
\]
respectively. Thus we obtain 
\[ \ov{F}(5,2)^5_3: \quad
\text{\setlength{\arraycolsep}{2.7pt}\scriptsize
$\displaystyle 
\begin{array}{c|cccc|cccc|cccc|cccc|cccc|}
 &\rt{12}&\rt{13}&\rt{14}&\rt{15}&\rt{21} &\rt{23}&\rt{24}&\rt{25}&\rt{31}&\rt{32} &\rt{34}&\rt{35}&\rt{41}&\rt{42}&\rt{43} &\rt{45} 
 &\rt{51}&\rt{52}&\rt{53} &\rt{54} \\ \hline 
 32 	&&&&			&&&&			&& &	&&& &	&		&&\cm&&\cm							\\ \hline
 42 	&&&&			&&&&\cm			&& && && &	&		&&\cm&&\cm				\\ 
 43 	&&&&			&&&&\cm			&& && && &	&		&&\cm&&\cm							\\ 
 \hline
\end{array}\; .$}
\]
Now let $j=2$. The initial colouring is
 \[ 
 \text{\setlength{\arraycolsep}{2.7pt}\scriptsize
$\displaystyle 
\begin{array}{c|cccc|cccc|cccc|cccc|cccc|}
 &\rt{12}&\rt{13}&\rt{14}&\rt{15}&\rt{21} &\rt{23}&\rt{24}&\rt{25}&\rt{31}&\rt{32} &\rt{34}&\rt{35}&\rt{41}&\rt{42}&\rt{43} &\rt{45} 
 &\rt{51}&\rt{52}&\rt{53} &\rt{54} \\ \hline 
 &0&0&0&0	
 &0&0&0&0	
 &0&0&&						 
 &0&0&&
 &0&0&& 	\\ \hline
\end{array}\;.$} 
\] 
In~\hyperlink{prop6.13_S1}{Step 1} of the algorithm, we 1-colour 54. In~\hyperlink{prop6.13_S2}{Step~ 2} this forces us
to 0-colour 53; this in turn forces us to 0-colour 43; this in turn
forces us to 0-colour 45; this in turn forces us to 0-colour 35; this
in turn forces us to 0-colour 34. Thus $I'_{2}(5,2)=\{54\}$.
Therefore we have
\[ 
\ov{F}(5,2)^5_2: \quad
\text{\setlength{\arraycolsep}{2.7pt}\scriptsize
$\displaystyle 
\begin{array}{c|cccc|cccc|cccc|cccc|cccc|}
 &\rt{12}&\rt{13}&\rt{14}&\rt{15}&\rt{21} &\rt{23}&\rt{24}&\rt{25}&\rt{31}&\rt{32} &\rt{34}&\rt{35}&\rt{41}&\rt{42}&\rt{43} &\rt{45} 
 &\rt{51}&\rt{52}&\rt{53} &\rt{54} \\ \hline 
 32 	&&&&			&&&&			&& &	&&& &	&		&&&&\cm							\\ \hline
 42 	&&&&			&&&&			&& && && &	&		&&&&\cm				\\ 
 43 	&&&&			&&&&			&& && && &	&		&&&&	\cm						\\ 
 \hline
\end{array}\; .$}
\]
We summarise this in the following table, where the last two blocks
are determined by $F(4,2)$ in Example~\ref{ex:free-patterns}\ref{exam6.12_i}:
\[
D(5,2): \quad
\text{\setlength{\arraycolsep}{2.7pt}\scriptsize
$\displaystyle 
\begin{array}{c|c||ccc||ccccc||ccccc|}
\hline j=&{2}	
&\multicolumn{3}{c||}{3}
&\multicolumn{5}{c||}{4}
&\multicolumn{5}{c|}{5}					\\ \hline 
 &\rt{54}&\rt{25}&\rt{52}&\rt{54}%
 &\rt{25}&\rt{32}&\rt{35}&\rt{52}%
 &\rt{53}&\rt{24 }&\rt{32 }&\rt{34 }%
 &\rt{42}&\rt{43}\\ 
 \hline
 32 &\cm &&\cm&\cm &&\cm&&\cm&\cm &&\cm&&\cm&\cm \\ 
 42 &\cm &\cm&\cm&\cm &\cm&\cm&\cm&\cm&\cm &\cm&\cm&\cm&\cm&\cm \\
 43 &\cm &\cm&\cm&\cm &\cm&\cm&\cm&\cm&\cm &\cm&\cm&\cm&\cm&\cm \\ \hline
\end{array}\,$}
\]
in which we omit columns that have no information, in order to save space.
We hence obtain 
\[
 F' (5,3): \quad
\text{\setlength{\arraycolsep}{2.7pt}\scriptsize
$\displaystyle 
\begin{array}{c|c|ccc|ccccc|ccccc|}
 &\rt{254}&\rt{325}&\rt{352}&\rt{354}%
 &\rt{425}&\rt{432}&\rt{435}&\rt{452}%
 &\rt{453}&\rt{524}&\rt{532}&\rt{534}%
 &\rt{542}&\rt{543}\\ \hline
 532 &\cm &&\cm&\cm &&\cm&&\cm&\cm &&\cm&&\cm&\cm \\ 
 542 &\cm &\cm&\cm&\cm &\cm&\cm&\cm&\cm&\cm &\cm&\cm&\cm&\cm&\cm \\
 543 &\cm &\cm&\cm&\cm &\cm&\cm&\cm&\cm&\cm &\cm&\cm&\cm&\cm&\cm \\ \hline
\end{array}\,.$}
\] 
\end{exam}




\begin{theo}\label{thm:properties}
 All the Properties~\ref{1}--\ref{4} hold for any $n \ge 2$, $r\ge
 1$.
\end{theo}

\begin{proof}
This follows from the results of this section by a double induction on
$n,r$. To be precise, we summarize the results proved in Propositions~\ref{prop:extend}, \ref{prop:decomp}, \ref{prop:free-ptns-exist}, and~\ref{prop:property-4} in the table below:
\[
\renewcommand{\arraystretch}{1.3}
\begin{tabular}{|c|l|c|}
\hline
 Result & \multicolumn{1}{c|}{Hypotheses} & Conclusion \\ \hline
\ref{prop:extend}& Prop.~1$(n-1,r)$, Prop.~2$(n,r-1)$ & Prop.~1$(n,r)$ \\
\ref{prop:decomp}& Prop.~1$(n-1,r)$, Prop.~4$(n-1,r-1)$ & Prop.~2$(n,r)$ \\
\ref{prop:free-ptns-exist}& Prop.~3$(n-1,r)$, Prop.~4$(n,r-1)$ & Prop.~3$(n,r)$ \\
\ref{prop:property-4}& Prop.~3$(n-1,r)$, Prop.~4$(n-1,r-1)$ & Prop.~4$(n,r)$\\ \hline
\end{tabular}\; .
\]
Properties~\ref{1} and~\ref{3} for $(n,1)$ are evident, for any
$n$. Properties~\ref{2} and~\ref{4} hold vacuously for $r=0$, for any
$n$, as there is nothing to do since $\E_\Bbbk(n,0) \cong
\Bbbk$. These serve as base cases for the induction.
\end{proof}
